\documentclass[10pt,a4paper]{amsart}
\usepackage{amsmath,amssymb,amsthm,graphicx}
% Retained preamble for the Stacks Project chapter snapshot in this directory.
% The chapter's own preamble.tex is a copy of the shared Stacks Project preamble; the block
% below is sliced MECHANICALLY, line for line, from it. Nothing was added, removed or
% redefined; the wrapper class, the xr-hyper/hyperref wiring and the externaldocument list
% are omitted because the assembler supplies the document class.
%
% lines taken from preamble.tex: 3, 5, 9, 11, 12, 156, 157, 158, ... (51 lines total)
%
% package coverage restored from the chapter's own preamble (2026-10-13): xy, xy-2cell

\usepackage[all]{xy}
\xyoption{2cell}
\UseAllTwocells

\newenvironment{reference}{\comment}{\endcomment}
\newenvironment{slogan}{\comment}{\endcomment}
\newenvironment{history}{\comment}{\endcomment}
\theoremstyle{plain}
\newtheorem{theorem}[subsection]{Theorem}
\newtheorem{proposition}[subsection]{Proposition}
\newtheorem{lemma}[subsection]{Lemma}
\theoremstyle{definition}
\newtheorem{definition}[subsection]{Definition}
\newtheorem{example}[subsection]{Example}
\newtheorem{exercise}[subsection]{Exercise}
\newtheorem{situation}[subsection]{Situation}
\theoremstyle{remark}
\newtheorem{remark}[subsection]{Remark}
\newtheorem{remarks}[subsection]{Remarks}
\numberwithin{equation}{subsection}
\def\lim{\mathop{\mathrm{lim}}\nolimits}
\def\colim{\mathop{\mathrm{colim}}\nolimits}
\def\Spec{\mathop{\mathrm{Spec}}}
\def\Hom{\mathop{\mathrm{Hom}}\nolimits}
\def\Ext{\mathop{\mathrm{Ext}}\nolimits}
\def\SheafHom{\mathop{\mathcal{H}\!\mathit{om}}\nolimits}
\def\SheafExt{\mathop{\mathcal{E}\!\mathit{xt}}\nolimits}
\def\Sch{\mathit{Sch}}
\def\Mor{\mathop{\mathrm{Mor}}\nolimits}
\def\Ob{\mathop{\mathrm{Ob}}\nolimits}
\def\Sh{\mathop{\mathit{Sh}}\nolimits}
\def\NL{\mathop{N\!L}\nolimits}
\def\CH{\mathop{\mathrm{CH}}\nolimits}
\def\proetale{{pro\text{-}\acute{e}tale}}
\def\etale{{\acute{e}tale}}
\def\QCoh{\mathit{QCoh}}
\def\Ker{\mathop{\mathrm{Ker}}}
\def\Im{\mathop{\mathrm{Im}}}
\def\Coker{\mathop{\mathrm{Coker}}}
\def\Coim{\mathop{\mathrm{Coim}}}
\DeclareMathSymbol{\boxtimes}{\mathbin}{AMSa}{"02}
\def\QCohstack{\mathcal{QC}\!\mathit{oh}}
\def\Cohstack{\mathcal{C}\!\mathit{oh}}
\def\Spacesstack{\mathcal{S}\!\mathit{paces}}
\def\Quotfunctor{\mathrm{Quot}}
\def\Hilbfunctor{\mathrm{Hilb}}
\def\Curvesstack{\mathcal{C}\!\mathit{urves}}
\def\Polarizedstack{\mathcal{P}\!\mathit{olarized}}
\def\Complexesstack{\mathcal{C}\!\mathit{omplexes}}
\def\Pic{\mathop{\mathrm{Pic}}\nolimits}
\def\Picardstack{\mathcal{P}\!\mathit{ic}}
\def\Picardfunctor{\mathrm{Pic}}
\def\Deformationcategory{\mathcal{D}\!\mathit{ef}}
\usepackage{cleveref}
\crefname{originallemma}{original lemma}{original lemmas}
\Crefname{originallemma}{Original lemma}{Original lemmas}
\makeatletter
\newlabel{lemma-quot-better-open}{{0DP3}{source}{}{sourcepointer-1}{}}
\newlabel{lemma-quot-better-open@cref}{{[originallemma][0][]0DP3}{[1][0][]source}{}{}{}}
\makeatother
\crefname{originallemma}{original lemma}{original lemmas}
\Crefname{originallemma}{Original lemma}{Original lemmas}
\makeatletter
\newlabel{lemma-quot-power-invertible}{{0DP8}{source}{}{sourcepointer-2}{}}
\newlabel{lemma-quot-power-invertible@cref}{{[originallemma][0][]0DP8}{[1][0][]source}{}{}{}}
\makeatother
\crefname{originallemma}{original lemma}{original lemmas}
\Crefname{originallemma}{Original lemma}{Original lemmas}
\makeatletter
\newlabel{lemma-quot-proper-over-base}{{0DPC}{source}{}{sourcepointer-3}{}}
\newlabel{lemma-quot-proper-over-base@cref}{{[originallemma][0][]0DPC}{[1][0][]source}{}{}{}}
\makeatother
\crefname{originallemma}{original lemma}{original lemmas}
\Crefname{originallemma}{Original lemma}{Original lemmas}
\makeatletter
\newlabel{lemma-quot-s-lfp}{{0DM3}{source}{}{sourcepointer-4}{}}
\newlabel{lemma-quot-s-lfp@cref}{{[originallemma][0][]0DM3}{[1][0][]source}{}{}{}}
\makeatother
\crefname{originallemma}{original lemma}{original lemmas}
\Crefname{originallemma}{Original lemma}{Original lemmas}
\makeatletter
\newlabel{morphisms-lemma-quasi-compact-immersion}{{01RG}{source}{}{sourcepointer-5}{}}
\newlabel{morphisms-lemma-quasi-compact-immersion@cref}{{[originallemma][0][]01RG}{[1][0][]source}{}{}{}}
\makeatother
\crefname{originallemma}{original lemma}{original lemmas}
\Crefname{originallemma}{Original lemma}{Original lemmas}
\makeatletter
\newlabel{properties-lemma-lift-finite-presentation}{{0G41}{source}{}{sourcepointer-6}{}}
\newlabel{properties-lemma-lift-finite-presentation@cref}{{[originallemma][0][]0G41}{[1][0][]source}{}{}{}}
\makeatother
\crefname{originallemma}{original lemma}{original lemmas}
\Crefname{originallemma}{Original lemma}{Original lemmas}
\makeatletter
\newlabel{remark-quot-numerical}{{0DP6}{source}{}{sourcepointer-7}{}}
\newlabel{remark-quot-numerical@cref}{{[originallemma][0][]0DP6}{[1][0][]source}{}{}{}}
\makeatother
\crefname{originallemma}{original lemma}{original lemmas}
\Crefname{originallemma}{Original lemma}{Original lemmas}
\makeatletter
\newlabel{spaces-divisors-definition-relatively-ample}{{0D31}{source}{}{sourcepointer-8}{}}
\newlabel{spaces-divisors-definition-relatively-ample@cref}{{[originallemma][0][]0D31}{[1][0][]source}{}{}{}}
\makeatother
\crefname{originallemma}{original lemma}{original lemmas}
\Crefname{originallemma}{Original lemma}{Original lemmas}
\makeatletter
\newlabel{spaces-limits-lemma-descend-finite-presentation}{{07SK}{source}{}{sourcepointer-9}{}}
\newlabel{spaces-limits-lemma-descend-finite-presentation@cref}{{[originallemma][0][]07SK}{[1][0][]source}{}{}{}}
\makeatother
\crefname{originallemma}{original lemma}{original lemmas}
\Crefname{originallemma}{Original lemma}{Original lemmas}
\makeatletter
\newlabel{spaces-limits-lemma-descend-invertible-modules}{{0D2X}{source}{}{sourcepointer-10}{}}
\newlabel{spaces-limits-lemma-descend-invertible-modules@cref}{{[originallemma][0][]0D2X}{[1][0][]source}{}{}{}}
\makeatother
\crefname{originallemma}{original lemma}{original lemmas}
\Crefname{originallemma}{Original lemma}{Original lemmas}
\makeatletter
\newlabel{spaces-limits-lemma-descend-modules-finite-presentation}{{07V7}{source}{}{sourcepointer-11}{}}
\newlabel{spaces-limits-lemma-descend-modules-finite-presentation@cref}{{[originallemma][0][]07V7}{[1][0][]source}{}{}{}}
\makeatother
\crefname{originallemma}{original lemma}{original lemmas}
\Crefname{originallemma}{Original lemma}{Original lemmas}
\makeatletter
\newlabel{spaces-limits-lemma-descend-separated}{{084T}{source}{}{sourcepointer-12}{}}
\newlabel{spaces-limits-lemma-descend-separated@cref}{{[originallemma][0][]084T}{[1][0][]source}{}{}{}}
\makeatother
\crefname{originallemma}{original lemma}{original lemmas}
\Crefname{originallemma}{Original lemma}{Original lemmas}
\makeatletter
\newlabel{spaces-morphisms-lemma-quasi-compact-local}{{03KG}{source}{}{sourcepointer-13}{}}
\newlabel{spaces-morphisms-lemma-quasi-compact-local@cref}{{[originallemma][0][]03KG}{[1][0][]source}{}{}{}}
\makeatother
\title{The Hilbert-polynomial piece of Quot is separated of finite presentation}
\author{The Stacks Project Authors}
\date{stacks-project:moduli.tex:local-snapshot-20261012; GNU FDL 1.2 or later}
\begin{document}
\maketitle
\noindent\textit{Editorial note.} Statement span is lines 1269-1281 of sources/stacks-moduli/pinned-original-source.tex; proof span is lines 1283-1339, the complete author proof (55 body lines).\par\medskip
\noindent\textit{Editorial note.} Navigation stubs: 13 cited-but-not-included labels are supplied as non-load-bearing pointers with their pinned tags where the snapshot tag index has them.\par\medskip
\noindent\textit{Editorial note.} extra\_packages is empty on purpose: the retained preamble already loads xy with options and a second declaration clashes in the assembler wrapper.\par\medskip
\noindent\textit{Editorial note.} Text is retained verbatim, including the author's own typos and placeholders.\par\medskip
\noindent\textit{Editorial note.} Queue index 4 of queue\_moduli.json (pb=2455, stmt 1269-1281, proof 1283-1339).\par\medskip
\noindent\textit{Editorial note.} stacks-moduli wave-1: fresh adjudication accept (2026-10-14); pool per Chief order 08:5x (dpa superseded to lane-s9); moduli.tex 53 labels/48 with proof judged 25 accept / 23 reject / 5 pb=0 env-dup rows excluded; spans cite-free and reference/slogan/history-free so no patchenv needed.\par\medskip
\noindent\textit{Allowed standard prerequisites (editorial).}\begin{itemize}
\item Explicit prerequisite boundary. The cited inputs are prerequisites; the argument in the retained span is not replaced by any cited result.
\end{itemize}\medskip
\ifdefined\hypertarget\hypertarget{sourcepointer-1}{}\fi
\noindent\textit{Source reference.} Original lemma 0DP3: sources/stacks-moduli/pinned-original-source.tex, source lines 870--881. Cited by the retained statement or proof but not part of the collected spans; supplied only as a navigation pointer. The argument does not depend on its content.\par\medskip
\ifdefined\hypertarget\hypertarget{sourcepointer-2}{}\fi
\noindent\textit{Source reference.} Original lemma 0DP8: sources/stacks-moduli/pinned-original-source.tex, source lines 1018--1027. Cited by the retained statement or proof but not part of the collected spans; supplied only as a navigation pointer. The argument does not depend on its content.\par\medskip
\ifdefined\hypertarget\hypertarget{sourcepointer-3}{}\fi
\noindent\textit{Source reference.} Original lemma 0DPC: sources/stacks-moduli/pinned-original-source.tex, source lines 1233--1244. Cited by the retained statement or proof but not part of the collected spans; supplied only as a navigation pointer. The argument does not depend on its content.\par\medskip
\ifdefined\hypertarget\hypertarget{sourcepointer-4}{}\fi
\noindent\textit{Source reference.} Original lemma 0DM3: sources/stacks-moduli/pinned-original-source.tex, source lines 755--759. Cited by the retained statement or proof but not part of the collected spans; supplied only as a navigation pointer. The argument does not depend on its content.\par\medskip
\ifdefined\hypertarget\hypertarget{sourcepointer-5}{}\fi
\noindent\textit{Source reference.} Original lemma 01RG: sources/stacks-moduli/pinned-original-source.tex, source lines 1339--1339. Cited by the retained statement or proof but not part of the collected spans; supplied only as a navigation pointer. The argument does not depend on its content.\par\medskip
\ifdefined\hypertarget\hypertarget{sourcepointer-6}{}\fi
\noindent\textit{Source reference.} Original lemma 0G41: sources/stacks-moduli/pinned-original-source.tex, source lines 1339--1339. Cited by the retained statement or proof but not part of the collected spans; supplied only as a navigation pointer. The argument does not depend on its content.\par\medskip
\ifdefined\hypertarget\hypertarget{sourcepointer-7}{}\fi
\noindent\textit{Source reference.} Original lemma 0DP6: sources/stacks-moduli/pinned-original-source.tex, source lines 950--981. Cited by the retained statement or proof but not part of the collected spans; supplied only as a navigation pointer. The argument does not depend on its content.\par\medskip
\ifdefined\hypertarget\hypertarget{sourcepointer-8}{}\fi
\noindent\textit{Source reference.} Original lemma 0D31: sources/stacks-moduli/pinned-original-source.tex, source lines 1339--1339. Cited by the retained statement or proof but not part of the collected spans; supplied only as a navigation pointer. The argument does not depend on its content.\par\medskip
\ifdefined\hypertarget\hypertarget{sourcepointer-9}{}\fi
\noindent\textit{Source reference.} Original lemma 07SK: sources/stacks-moduli/pinned-original-source.tex, source lines 1339--1339. Cited by the retained statement or proof but not part of the collected spans; supplied only as a navigation pointer. The argument does not depend on its content.\par\medskip
\ifdefined\hypertarget\hypertarget{sourcepointer-10}{}\fi
\noindent\textit{Source reference.} Original lemma 0D2X: sources/stacks-moduli/pinned-original-source.tex, source lines 1339--1339. Cited by the retained statement or proof but not part of the collected spans; supplied only as a navigation pointer. The argument does not depend on its content.\par\medskip
\ifdefined\hypertarget\hypertarget{sourcepointer-11}{}\fi
\noindent\textit{Source reference.} Original lemma 07V7: sources/stacks-moduli/pinned-original-source.tex, source lines 1339--1339. Cited by the retained statement or proof but not part of the collected spans; supplied only as a navigation pointer. The argument does not depend on its content.\par\medskip
\ifdefined\hypertarget\hypertarget{sourcepointer-12}{}\fi
\noindent\textit{Source reference.} Original lemma 084T: sources/stacks-moduli/pinned-original-source.tex, source lines 1339--1339. Cited by the retained statement or proof but not part of the collected spans; supplied only as a navigation pointer. The argument does not depend on its content.\par\medskip
\ifdefined\hypertarget\hypertarget{sourcepointer-13}{}\fi
\noindent\textit{Source reference.} Original lemma 03KG: sources/stacks-moduli/pinned-original-source.tex, source lines 1339--1339. Cited by the retained statement or proof but not part of the collected spans; supplied only as a navigation pointer. The argument does not depend on its content.\par\medskip
\begin{lemma}
\label{lemma-quot-qc-over-base}
Let $f : X \to B$ be a separated morphism of finite presentation
of algebraic spaces. Let $\mathcal{F}$ be a finitely presented
$\mathcal{O}_X$-module. Let $\mathcal{L}$ be an invertible
$\mathcal{O}_X$-module ample on $X/B$, see
Divisors on Spaces, Definition
\ref{spaces-divisors-definition-relatively-ample}.
The algebraic space $\Quotfunctor^P_{\mathcal{F}/X/B}$
parametrizing quotients of $\mathcal{F}$
having Hilbert polynomial $P$ with respect to $\mathcal{L}$
is separated of finite presentation over $B$.
\end{lemma}

\begin{proof}
We have already seen that $\Quotfunctor_{\mathcal{F}/X/B} \to B$
is separated and locally of finite presentation, see
Lemma \ref{lemma-quot-s-lfp}. Thus it suffices to show that
the open subspace $\Quotfunctor^P_{\mathcal{F}/X/B}$
of Remark \ref{remark-quot-numerical} is quasi-compact over $B$.

\medskip\noindent
The question is \'etale local on $B$
(Morphisms of Spaces, Lemma \ref{spaces-morphisms-lemma-quasi-compact-local}).
Thus we may assume $B$ is affine.

\medskip\noindent
Assume $B = \Spec(\Lambda)$. Write $\Lambda = \colim \Lambda_i$ as the
colimit of its finite type $\mathbf{Z}$-subalgebras. Then
we can find an $i$ and a system $X_i, \mathcal{F}_i, \mathcal{L}_i$
as in the lemma over $B_i = \Spec(\Lambda_i)$ whose base change to
$B$ gives $X, \mathcal{F}, \mathcal{L}$.
This follows from
Limits of Spaces, Lemmas
\ref{spaces-limits-lemma-descend-finite-presentation} (to find $X_i$),
\ref{spaces-limits-lemma-descend-modules-finite-presentation} (to find
$\mathcal{F}_i$), \ref{spaces-limits-lemma-descend-invertible-modules}
(to find $\mathcal{L}_i$), and \ref{spaces-limits-lemma-descend-separated}
(to make $X_i$ separated). Because
$$
\Quotfunctor_{\mathcal{F}/X/B} = B \times_{B_i}
\Quotfunctor_{\mathcal{F}_i/X_i/B_i}
$$
and similarly for $\Quotfunctor^P_{\mathcal{F}/X/B}$ we reduce
to the case discussed in the next paragraph.

\medskip\noindent
Assume $B$ is affine and Noetherian. We may replace $\mathcal{L}$
by a positive power, see Lemma \ref{lemma-quot-power-invertible}.
Thus we may assume there exists an immersion $i : X \to \mathbf{P}^n_B$
such that $i^*\mathcal{O}_{\mathbf{P}^n}(1) = \mathcal{L}$. By
Morphisms, Lemma \ref{morphisms-lemma-quasi-compact-immersion}
there exists a closed subscheme $X' \subset \mathbf{P}^n_B$
such that $i$ factors through an open immersion $j : X \to X'$.
By Properties, Lemma \ref{properties-lemma-lift-finite-presentation}
there exists a finitely presented $\mathcal{O}_{X'}$-module
$\mathcal{G}$ such that $j^*\mathcal{G} = \mathcal{F}$.
Thus we obtain an open immersion
$$
\Quotfunctor_{\mathcal{F}/X/B}
\longrightarrow
\Quotfunctor_{\mathcal{G}/X'/B}
$$
by Lemma \ref{lemma-quot-better-open}. Clearly this open immersion
sends $\Quotfunctor^P_{\mathcal{F}/X/B}$ into
$\Quotfunctor^P_{\mathcal{G}/X'/B}$. Now
$\Quotfunctor^P_{\mathcal{G}/X'/B}$ is proper over $B$ by
Lemma \ref{lemma-quot-proper-over-base}.
Therefore it is Noetherian and since any open of a Noetherian
algebraic space is quasi-compact we win.
\end{proof}

\end{document}
